Emphasize that not all critical transitions are predictable under the currently developed framework and cite criticism within the community.
Stress how the current theory is largely developed for stochastically forced systems that experience a slow passage through a saddle-node bifurcation.
The goal of this section is to select a candidate timeseries from the real world (given in Motivation) and fit a saddle-node model with stochastic fluctuations.
To do so we need two main ingredients: geometric singular perturbation theory (or Fenichel theory) and large deviation theory (or Freidlin and Wentzell theory).

\section{Nonautonomous dynamics}\label{subsec:nonautonomous}
Introduce the autonomous saddle-node normal form and an example of coordinate transformations for the saddle-node bifurcation theorem.
Extend it to the dynamic case with the notion of pullback attractors.

\section{Stochastic calculus}\label{subsec:stochastic}
Introduce the fundamentals of stochastic calculus including the OUP and the FPE.

\section{Large deviations in slow-fast systems}\label{subsec:slowfast}
State that the correct mathematical interpretation of critical transitions is that of rare events in small perturbations of dynamic saddle-node systems.
Introduce the two ingredients separately.

\subsection{Geometric singular perturbation theory}\label{subsubsec:gspt}
Main results of Fenichel theory.

\subsection{Large deviation theory}\label{subsubsec:ldt}
Main results of Freidlin-Wentzell theory.


